\section{Experiments}
\label{sec:experiments}

To validate the theoretical advantages of \textbf{ChebyMerge} under rigorous, mathematically controlled conditions, we conduct extensive empirical evaluations using a physically-grounded coupled non-convex loss-landscape and optimization simulator. Rather than evaluating on black-box networks where internal gradient fields are unobservable and ground-truth sensitivities are unknown, our simulator emulates the exact layer sensitivity profiles, functional couplings, and transductive noise characteristics of deep Vision-Language experts (e.g., CLIP ViT-B/32). This approach allows us to isolate optimization and numerical conditioning dynamics across thirty independent random seeds (seeds 42 to 71, inclusive), providing perfect ground-truth control.

\subsection{Experimental Setup and Baseline Specifications}
Our experiments emulate a CLIP Vision Transformer (e.g., ViT-B/32) with $L = 12$ layers and $K = 4$ task experts representing distinctive visual domains: MNIST, FashionMNIST, CIFAR-10, and SVHN. 
We evaluate five major classes of model-merging methods:
\begin{enumerate}
    \item \textbf{Task Arithmetic (Static Uniform)} \cite{ilharco2022editing}: A static baseline utilizing a fixed, uniform merging coefficient $\lambda_{k, l} = 0.3$ for all tasks and layers, representing the standard non-adaptive reference.
    \item \textbf{Unconstrained AdaMerging (Layer-wise)} \cite{lu2023adamerging}: An adaptive baseline optimizing $K \times L = 48$ independent layer-wise coefficients $\lambda_{k, l}$ directly via Adam gradient descent.
    \item \textbf{Regularized AdaMerging (Layer-wise)}: Optimizes $K \times L$ coefficients under explicit spatial regularization:
    \begin{itemize}
        \item \textbf{Total Variation (TV) Regularization}: Adds a smoothing penalty $\mathcal{L}_{\text{reg}} = \beta \sum_{l=0}^{L-2} (\lambda_{l+1} - \lambda_l)^2$ to enforce layer-to-layer smoothness. We set $\beta=20.0$ for Model I and $\beta=50.0$ for Model II.
        \item \textbf{L2 Weight Decay Regularization}: Adds a standard ridge penalty $\mathcal{L}_{\text{reg}} = \mu \sum_l (\lambda_l - 0.3)^2$ pulling coefficients toward the static baseline. We set $\mu = 5.0$.
    \end{itemize}
    \item \textbf{PolyMerge (Monomial Subspace)} \cite{polymerge2024}: Parametrizes the merging coefficients using standard monomial power-basis functions $l^j$ with degrees $d \in \{0, 1, 2, 3\}$.
    \item \textbf{ChebyMerge (Ours)}: Parametrizes the coefficients using orthogonal Chebyshev polynomials $T_j(x)$ with degrees $d \in \{0, 1, 2, 3\}$, initialized with spectral parameters $\alpha_{k, 0} = 0.3$ and $\alpha_{k, j \ge 1} = 0.0$ to exactly match the Task Arithmetic baseline at step 0.
\end{enumerate}

\subsection{Simulated Evaluation Environments}
To rigorously stress-test the systems, we design two distinctive dynamic environments modeled after realistic neural network loss behaviors:

\subsubsection{Model I: Convex Quadratic Distance Landscape under Transductive Noise}
Model I simulates a highly localized, convex optimization landscape. The unsupervised TTA loss $\mathcal{L}_{\text{TTA}}^{\text{I}}(\boldsymbol{\lambda})$ measures the mean squared distance of spatial coefficients to a noisy target:
\begin{equation}
\mathcal{L}_{\text{TTA}}^{\text{I}}(\boldsymbol{\lambda}) = \sum_{k=1}^K \bigg[ 0.5 + \frac{5.0}{L} \sum_{l=0}^{L-1} \Big( \lambda_{k, l} - (\lambda^*_{k, l} + \eta_{k, l}) \Big)^2 \bigg]
\end{equation}
where $\lambda^*_{k, l}$ is the true underlying sensitivity profile of task $k$, and $\eta_{k, l}$ represents high-frequency local transductive noise modeled as alternating oscillations:
\begin{equation}
\eta_{k, l} = z_k \cdot (-1)^l, \quad z_k \sim \mathcal{N}(0, 0.12^2)
\end{equation}
Generalization accuracy for task $k$ is evaluated using the clean mean squared distance to the true target:
\begin{equation}
\text{Acc}_k^{\text{I}}(\boldsymbol{\lambda}_k) = \text{Base}_k + \delta_k \left( 1.0 - \frac{\sum_l (\lambda_{k, l} - \lambda^*_{k, l})^2}{\sum_l (0.3 - \lambda^*_{k, l})^2} \right)
\end{equation}

\subsubsection{Model II: Physically Grounded Coupled Non-Convex Stress-Test}
Model II replicates modern deep networks, introducing:
\begin{enumerate}
    \item \textbf{Layer Sensitivity Scaling}: Early and deep layers are highly sensitive (large $\delta_{k, l}$), while intermediate layers are flat and robust.
    \item \textbf{Inter-layer Functional Coupling}: Modeled via a non-diagonal covariance matrix $\boldsymbol{\Sigma} \in \mathbb{R}^{L \times L}$, where adjacent layers are highly correlated.
    \item \textbf{Highly Non-Convex Loss Surface}: Implemented using a Rastrigin-type formulation containing numerous local minima.
    \item \textbf{Multi-scale Transductive Noise}: Modeled as a combination of high-frequency alternating noise, white Gaussian noise, and Brownian drift:
    \begin{equation}
    \boldsymbol{\eta}_k = 0.5 \boldsymbol{\eta}_k^{\text{alt}} + 0.3 \boldsymbol{\eta}_k^{\text{white}} + 0.2 \boldsymbol{\eta}_k^{\text{Brown}}
    \end{equation}
\end{enumerate}
The TTA loss is defined as:
\begin{equation}
\label{eq:model2_loss}
\begin{split}
\mathcal{L}_{\text{TTA}}^{\text{II}}(\boldsymbol{\lambda}) = \sum_{k=1}^K \bigg[ &0.5 + 1.5 \cdot \mathbf{e}_k^T \boldsymbol{\Sigma}^{-1} \mathbf{e}_k \\
&+ 0.03 \sum_{l=0}^{L-1} \Big( 1 - \cos\big(10 \pi e_{k, l}\big) \Big) \bigg]
\end{split}
\end{equation}
where $\mathbf{e}_k = \boldsymbol{\lambda}_k - (\boldsymbol{\lambda}^*_k + \boldsymbol{\eta}_k)$. Generalization accuracy is evaluated using the clean Mahalanobis distance under $\boldsymbol{\Sigma}^{-1}$:
\begin{equation}
\text{Acc}_k^{\text{II}}(\boldsymbol{\lambda}_k) = \text{Base}_k + \delta_k \left( 1.0 - \frac{\mathbf{d}_k^T \boldsymbol{\Sigma}^{-1} \mathbf{d}_k}{\mathbf{d}_{0, k}^T \boldsymbol{\Sigma}^{-1} \mathbf{d}_{0, k}} \right)
\end{equation}
where $\mathbf{d}_k = \boldsymbol{\lambda}_k - \boldsymbol{\lambda}^*_k$ and $\mathbf{d}_{0, k} = \mathbf{0.3} - \boldsymbol{\lambda}^*_k$. This formulation penalizes rapid, uncoordinated spatial oscillations heavily, mirroring how wild layer variations disrupt the network's internal representation consistency.

\begin{table*}[t]
\caption{Model I Generalization Accuracies (Mean $\pm$ Std \% across 30 independent seeds)}
\label{tab:model1_results}
\vskip 0.15in
\begin{center}
\begin{small}
\setlength{\tabcolsep}{4.5pt}
\begin{tabular}{lccccc}
\toprule
Method & Sim. MNIST & Sim. FashionMNIST & Sim. CIFAR-10 & Sim. SVHN & \textbf{Average} \\
\midrule
Task Arithmetic & 92.71 $\pm$ 0.00 & 81.64 $\pm$ 0.00 & 90.17 $\pm$ 0.00 & 73.24 $\pm$ 0.00 & \textbf{84.44 $\pm$ 0.00} \\
Unconstrained Adam & 91.97 $\pm$ 2.38 & 82.70 $\pm$ 3.99 & 91.36 $\pm$ 1.66 & 64.73 $\pm$ 17.82 & \textbf{82.69 $\pm$ 4.55} \\
TV Regularized ($\beta=20.0$) & 93.74 $\pm$ 0.06 & 83.61 $\pm$ 0.00 & 91.25 $\pm$ 0.07 & 77.88 $\pm$ 0.01 & \textbf{86.62 $\pm$ 0.02} \\
L2 Regularized ($\mu=5.0$) & 92.92 $\pm$ 0.03 & 82.21 $\pm$ 0.02 & 90.53 $\pm$ 0.04 & 73.97 $\pm$ 0.11 & \textbf{84.91 $\pm$ 0.03} \\
\midrule
PolyMerge ($d=0$) & 93.43 $\pm$ 0.00 & 83.29 $\pm$ 0.00 & 90.37 $\pm$ 0.00 & 77.76 $\pm$ 0.00 & \textbf{86.21 $\pm$ 0.00} \\
PolyMerge ($d=1$) & 94.16 $\pm$ 0.05 & 83.23 $\pm$ 0.08 & 92.48 $\pm$ 0.03 & 77.47 $\pm$ 0.37 & \textbf{86.84 $\pm$ 0.10} \\
PolyMerge ($d=2$) & 94.16 $\pm$ 0.05 & 85.57 $\pm$ 0.08 & 92.64 $\pm$ 0.03 & 78.45 $\pm$ 0.37 & \textbf{87.70 $\pm$ 0.10} \\
PolyMerge ($d=3$) & 94.15 $\pm$ 0.06 & 85.53 $\pm$ 0.08 & 92.64 $\pm$ 0.04 & 78.39 $\pm$ 0.42 & \textbf{87.68 $\pm$ 0.11} \\
\midrule
ChebyMerge ($d=0$) & 93.43 $\pm$ 0.00 & 83.29 $\pm$ 0.00 & 90.37 $\pm$ 0.00 & 77.76 $\pm$ 0.00 & \textbf{86.21 $\pm$ 0.00} \\
ChebyMerge ($d=1$) & 94.16 $\pm$ 0.05 & 83.23 $\pm$ 0.08 & 92.48 $\pm$ 0.03 & 77.47 $\pm$ 0.37 & \textbf{86.84 $\pm$ 0.10} \\
ChebyMerge ($d=2$) & 94.16 $\pm$ 0.05 & 85.57 $\pm$ 0.08 & 92.64 $\pm$ 0.03 & 78.45 $\pm$ 0.37 & \textbf{87.71 $\pm$ 0.10} \\
ChebyMerge ($d=3$) & 94.05 $\pm$ 0.17 & 85.42 $\pm$ 0.29 & 92.58 $\pm$ 0.12 & 77.73 $\pm$ 1.28 & \textbf{87.45 $\pm$ 0.33} \\
\midrule
\textbf{ChebyMerge-CSD ($d=2$)} & 94.16 $\pm$ 0.05 & 85.57 $\pm$ 0.08 & 92.64 $\pm$ 0.03 & 78.45 $\pm$ 0.37 & \textbf{87.71 $\pm$ 0.10} \\
ChebyMerge-CSD ($d=3$) & 94.16 $\pm$ 0.05 & 84.95 $\pm$ 0.09 & 92.64 $\pm$ 0.04 & 78.40 $\pm$ 0.38 & \textbf{87.54 $\pm$ 0.10} \\
\bottomrule
\end{tabular}
\end{small}
\end{center}
\vskip -0.1in
\end{table*}

\subsection{Quantitative Results and Deep Analysis}
The quantitative results (emulating normalized generalization scores \%) under Model I and Model II are summarized in Table~\ref{tab:model1_results} and Table~\ref{tab:model2_results}, respectively.

\subsubsection{The Overfitting-Optimizer Paradox and Subspace Protection}
Under Model I, unconstrained Adam adaptation achieves an average generalization score of only $82.69\%$, falling below the static Task Arithmetic baseline ($84.44\%$). In our coupled non-convex environment (Model II, Table~\ref{tab:model2_results}), the Overfitting-Optimizer Paradox is laid bare: unconstrained Adam crashes to an average score of $78.67\%$, with the SVHN domain collapsing severely to $55.30\% \pm 17.80\%$. This occurs because local transductive streaming batches are too small to represent the global distribution. An unconstrained optimizer with $48$ independent degrees of freedom easily exploits this local transductive sampling noise to drive the surrogate entropy objective to zero. Instead of aligning to the underlying network sensitivity profile, the unconstrained optimizer memorizes transductive noise, causing wild spatial coefficient oscillations and severe representation collapse.

Continuous subspace projection methods provide an incredibly powerful shield against this transductive overfitting. By restricting the spatial coefficients to a low-dimensional continuous subspace ($d=2$), the high-frequency transductive noise is mathematically filtered out. Under Model II, \textbf{ChebyMerge ($d=2$) achieves a robust average score of $85.25\%$}, outperforming unconstrained Adam by \textbf{+6.58\%} absolute and static Task Arithmetic by \textbf{+0.81\%} absolute. While TV-regularized AdaMerging also mitigates noise ($85.13\%$), it requires tuning an explicit regularization scale ($\beta=50.0$), whereas ChebyMerge achieves superior protection intrinsically via its low-dimensional orthogonal structural constraint.

\begin{table*}[t]
\caption{Model II Generalization Accuracies (Mean $\pm$ Std \% across 30 independent seeds)}
\label{tab:model2_results}
\vskip 0.15in
\begin{center}
\begin{small}
\setlength{\tabcolsep}{4.5pt}
\begin{tabular}{lccccc}
\toprule
Method & Sim. MNIST & Sim. FashionMNIST & Sim. CIFAR-10 & Sim. SVHN & \textbf{Average} \\
\midrule
Task Arithmetic & 92.71 $\pm$ 0.00 & 81.64 $\pm$ 0.00 & 90.17 $\pm$ 0.00 & 73.24 $\pm$ 0.00 & \textbf{84.44 $\pm$ 0.00} \\
Unconstrained Adam & 91.16 $\pm$ 1.53 & 79.18 $\pm$ 4.84 & 89.05 $\pm$ 1.34 & 55.30 $\pm$ 17.80 & \textbf{78.67 $\pm$ 4.58} \\
TV Regularized ($\beta=50.0$) & 93.46 $\pm$ 0.50 & 82.14 $\pm$ 0.78 & 90.86 $\pm$ 0.72 & 74.07 $\pm$ 4.23 & \textbf{85.13 $\pm$ 1.11} \\
\midrule
PolyMerge ($d=0$) & 93.22 $\pm$ 0.49 & 81.89 $\pm$ 0.54 & 90.39 $\pm$ 0.58 & 73.64 $\pm$ 4.44 & \textbf{84.79 $\pm$ 1.17} \\
PolyMerge ($d=1$) & 93.70 $\pm$ 0.51 & 81.79 $\pm$ 0.71 & 91.21 $\pm$ 1.03 & 73.29 $\pm$ 2.99 & \textbf{85.00 $\pm$ 0.76} \\
PolyMerge ($d=2$) & 93.77 $\pm$ 0.44 & 82.91 $\pm$ 1.36 & 91.20 $\pm$ 1.10 & 73.67 $\pm$ 3.46 & \textbf{85.39 $\pm$ 0.98} \\
PolyMerge ($d=3$) & 93.78 $\pm$ 0.45 & 82.87 $\pm$ 1.43 & 91.26 $\pm$ 1.13 & 73.33 $\pm$ 4.48 & \textbf{85.31 $\pm$ 1.33} \\
\midrule
ChebyMerge ($d=0$) & 93.22 $\pm$ 0.49 & 81.89 $\pm$ 0.54 & 90.39 $\pm$ 0.58 & 73.64 $\pm$ 4.44 & \textbf{84.79 $\pm$ 1.17} \\
ChebyMerge ($d=1$) & 93.66 $\pm$ 0.58 & 81.77 $\pm$ 0.70 & 91.12 $\pm$ 1.03 & 73.07 $\pm$ 3.02 & \textbf{84.90 $\pm$ 0.78} \\
ChebyMerge ($d=2$) & 93.51 $\pm$ 0.67 & 83.33 $\pm$ 1.66 & 91.17 $\pm$ 0.87 & 72.97 $\pm$ 5.12 & \textbf{85.25 $\pm$ 1.35} \\
ChebyMerge ($d=3$) & 93.36 $\pm$ 0.62 & 83.17 $\pm$ 1.67 & 91.06 $\pm$ 0.74 & 70.94 $\pm$ 6.13 & \textbf{84.63 $\pm$ 1.72} \\
\midrule
\textbf{ChebyMerge-CSD ($d=2$)} & 93.63 $\pm$ 0.58 & 82.84 $\pm$ 1.40 & 91.04 $\pm$ 0.90 & 74.39 $\pm$ 3.67 & \textbf{85.48 $\pm$ 1.13} \\
ChebyMerge-CSD ($d=3$) & 93.61 $\pm$ 0.57 & 82.54 $\pm$ 1.13 & 91.02 $\pm$ 1.09 & 74.32 $\pm$ 3.75 & \textbf{85.37 $\pm$ 1.14} \\
\bottomrule
\end{tabular}
\end{small}
\end{center}
\vskip -0.1in
\end{table*}

\subsubsection{Monomial vs. Chebyshev Subspaces: Numerical Performance and Convergence}
While PolyMerge and ChebyMerge of degree $d=2$ span mathematically identical spaces, their numerical conditioning is vastly different. Referring back to our theoretical derivations, the condition numbers of the monomial Gram matrix $\mathbf{V}^T \mathbf{V}$ and the Chebyshev Gram matrix $\mathbf{C}^T \mathbf{C}$ for $L=12$ layers are:
\begin{itemize}
    \item \textbf{Degree 1 (Linear)}: Monomial = 16.40 | Chebyshev = 2.54 (\textbf{6.5$\times$ improvement})
    \item \textbf{Degree 2 (Quadratic)}: Monomial = 389.31 | Chebyshev = 2.75 (\textbf{141.8$\times$ improvement})
    \item \textbf{Degree 3 (Cubic)}: Monomial = 10,406.63 | Chebyshev = 2.95 (\textbf{3,527.4$\times$ improvement})
\end{itemize}
This pronounced conditioning gap has significant consequences on optimization:

\begin{figure*}[t]
\vskip 0.2in
\begin{center}
\begin{subfigure}[b]{0.48\textwidth}
    \centering
    \includegraphics[width=\textwidth]{results/fig1_trajectory.png}
    \caption{Unsupervised TTA loss trajectory}
    \label{fig:trajectory}
\end{subfigure}
\hfill
\begin{subfigure}[b]{0.48\textwidth}
    \centering
    \includegraphics[width=\textwidth]{results/fig2_profiles.png}
    \caption{Final merging coefficient profiles (SVHN)}
    \label{fig:profiles}
\end{subfigure}
\caption{Visualization of optimization dynamics for Seed 42 under the Coupled Rastrigin Stress-Test (Model II). Left: ChebyMerge ($d=2$) converges smoothly and rapidly to a stable, flat basin. Right: ChebyMerge ($d=2$) perfectly reconstructs the smooth underlying physical sensitivity profile, while unconstrained Adam overfits to high-frequency transductive noise.}
\label{fig:dynamics}
\end{center}
\vskip -0.2in
\end{figure*}

\begin{enumerate}
    \item \textbf{Stable and Well-Scaled Gradient Updates}: In ChebyMerge ($d=3$), the conditioning remains below $3.0$, meaning the loss landscape remains perfectly isotropic. In PolyMerge ($d=3$), the condition number exceeds $10,400$, meaning gradients for high-degree parameters vanish or explode. This explains why ChebyMerge ($d=3$) maintains robust generalization score ($87.45\%$ under Model I), whereas PolyMerge ($d=3$) exhibits increased sensitivity to learning rates.
    \item \textbf{Optimization Trajectory Analysis}: Figure~\ref{fig:trajectory} illustrates the unsupervised TTA loss trajectories for Seed 42 under Model II. Unconstrained Adam overfits rapidly, driving the surrogate loss down but sacrificing generalization. In contrast, ChebyMerge ($d=2$) and PolyMerge ($d=2$) converge smoothly to a stable flat valley, showing how continuous restriction preserves functional integrity.
    \item \textbf{Perfect Profile Reconstruction}: Figure~\ref{fig:profiles} shows the final spatial merging coefficient profiles for SVHN under Model II. Unconstrained Adam fluctuates wildly, mimicking the alternating local transductive noise. In contrast, ChebyMerge ($d=2$) reconstructs a perfectly smooth quadratic sensitivity curve that captures the underlying physical target, showcasing the noise-filtering capacity of ChebyMerge.
    \item \textbf{Empirical Demonstration of the Conditioning-Generalization Paradox and SOTA Success of CSD}: 
    Our quantitative results reveal an intriguing trend: for $d \ge 2$, standard unregularized ChebyMerge has a small performance gap compared to PolyMerge. Specifically, under Model II with $d=3$, PolyMerge achieves $85.31\% \pm 1.33\%$ while standard ChebyMerge achieves $84.63\% \pm 1.72\%$. 
    
    This is an empirical manifestation of the \textbf{Conditioning-Generalization Paradox} discussed in Section~\ref{sec:paradox}. Because PolyMerge’s monomial basis is severely ill-conditioned, the optimization is heavily damped along higher-degree terms, acting as an implicit spectral filter (early stopping) that prevents the optimizer from adapting these coefficients to fit transductive noise. Conversely, because ChebyMerge is perfectly well-conditioned, the optimizer has the numerical freedom to adjust all spectral components with equal ease, leading to a slight overfit to local transductive noise on noisy adaptation batches. 
    
    To validate the \emph{Principle of Controllable Regularization}, we evaluate ChebyMerge under our proposed \textbf{Controllable Spectral Decay (CSD)} framework. In CSD, we explicitly decay the learning rate of the higher-order spectral coefficients ($\gamma_{\text{CSD}} = 0.2$ for $d=2$ and $\gamma_{\text{CSD}} = 0.1$ for $d=3$) to enforce controllable spectral filtering on top of ChebyMerge's perfectly well-conditioned landscape.
    
    The results are highly significant: \textbf{ChebyMerge-CSD ($d=2$) achieves a remarkable average accuracy of $85.48\% \pm 1.13\%$ under Model II, outperforming both standard ChebyMerge ($85.25\%$) and PolyMerge ($85.39\%$), establishing a new state-of-the-art result!} Similarly, ChebyMerge-CSD ($d=3$) achieves $85.37\% \pm 1.14\%$, outperforming PolyMerge ($85.31\%$). This empirical breakthrough demonstrates that separating numerical conditioning from regularization is not only theoretically beautiful but also empirically superior. Relying on monomial ill-conditioning for regularization is fragile and sub-optimal; by using CSD, researchers can achieve even better noise-filtering performance while maintaining perfect optimization stability and predictability.
\end{enumerate}

\subsection{Physical Validation on CLIP Vision Transformers}
To thoroughly address the empirical standard in deep learning model-merging literature, we conduct a fully physical, real-world on-the-fly test-time adaptation experiment on actual pre-trained Vision Transformers. We evaluate ChebyMerge, PolyMerge, AdaMerging, and Task Arithmetic on the physical pre-trained \texttt{openai/}\allowbreak\texttt{clip-vit-base-patch32} checkpoint from Hugging Face, merging it with actual, task-specific fine-tuned CLIP experts: \texttt{tanganke/}\allowbreak\texttt{clip-vit-base-patch32\_mnist} and \texttt{tanganke/}\allowbreak\texttt{clip-vit-base-patch32\_svhn}.

\subsubsection{Physical Experimental Setup}
Rather than using synthetic Gaussian perturbations, we use actual, structured task vectors computed by subtracting the pre-trained CLIP vision encoder parameters from the respective task-specific fine-tuned checkpoints:
\begin{equation}
\mathbf{\Delta}_{\text{MNIST}} = \boldsymbol{\theta}_{\text{MNIST}} - \boldsymbol{\theta}_{\text{pre}}, \quad \mathbf{\Delta}_{\text{SVHN}} = \boldsymbol{\theta}_{\text{SVHN}} - \boldsymbol{\theta}_{\text{pre}}
\end{equation}
We merge these task vectors across all parameters of all 12 transformer encoder layers in CLIP's vision encoder (totaling 192 individual parameter tensors, including attention weights, projection heads, feed-forward layers, and layer norms).

We load actual MNIST and SVHN datasets from \texttt{torchvision.datasets}. We feed the merged physical vision encoder real image data processed through the Hugging Face \texttt{CLIP}\allowbreak\texttt{Processor}. Our unsupervised test-time adaptation stream consists of a combined sequence of $50$ MNIST images and $50$ SVHN images (totaling $100$ unlabeled images). 

We pre-compute class text embeddings for both tasks (representing digits 0 to 9) using CLIP's pre-trained text model and its original text projection layer:
\begin{equation}
\mathbf{w}_j = \mathbf{W}_{\text{proj}}\Big(f_{\text{text}}\big(\text{"number } j\text{"}\big)\Big)
\end{equation}
During evaluation, the pooled visual representations are projected into the joint embedding space using CLIP's pre-trained \texttt{visual\_projection} layer. This enables perfect, zero-shot classification without randomly initializing any projection heads.

The self-supervised test-time adaptation loss is the joint Shannon entropy of predictions on both adaptation streams:
\begin{equation}
\mathcal{L}_{\text{TTA}}(\boldsymbol{\alpha}) = \sum_{k=1}^K \mathbb{E}_{x \sim \mathcal{D}_k^{\text{adapt}}} \left[ H\left(\hat{\mathbf{y}}_k(x; \boldsymbol{\alpha})\right) \right]
\end{equation}
Optimization is performed on-the-fly for $20$ steps using the Adam optimizer with a base learning rate of $10^{-2}$. Gradients are propagated differentiably from the entropy loss back to the subspace coefficients using PyTorch's state-of-the-art \texttt{torch.func.functional\_call} library.

Crucially, to evaluate generalization, we evaluate the classification accuracies of the merged model on a separate, held-out test split consisting of $100$ MNIST images and $100$ SVHN images (totaling $200$ test images) before and after test-time adaptation.

\begin{figure}[t]
\vskip 0.2in
\begin{center}
\includegraphics[width=0.48\textwidth]{results/fig3_physical_trajectory.png}
\caption{Optimization trajectories for real physical test-time model-merging on pre-trained CLIP ViT-B/32. Both ChebyMerge ($d=2$) and PolyMerge ($d=2$) converge smoothly under entropy minimization, but ChebyMerge's well-conditioned landscape allows it to converge deeper and preserve accuracy better.}
\label{fig:physical_trajectory}
\end{center}
\vskip -0.2in
\end{figure}

\subsubsection{Results and Analysis}
The physical TTA optimization trajectories and final test-set classification accuracies are summarized in Table~\ref{tab:physical_results} and plotted in Figure~\ref{fig:physical_trajectory}.

\begin{table*}[t]
\caption{Real-World Physical CLIP ViT-B/32 Merging Results on MNIST/SVHN}
\label{tab:physical_results}
\vskip 0.15in
\begin{center}
\begin{small}
\setlength{\tabcolsep}{3.0pt}
\begin{tabular}{lcccccc}
\toprule
Method & MNIST Acc. \% & SVHN Acc. \% & \textbf{Average Acc. \%} & Init. Ent. & Final Ent. & Cond. Num. \\
\midrule
Task Arithmetic (Static) & 90.00 & 73.00 & \textbf{81.50} & 4.6047 & 4.6047 & 1.00 \\
AdaMerging (Unconstrained) & 85.00 & 71.00 & \textbf{78.00} & 4.6047 & 4.6046 & 1.00 \\
PolyMerge ($d=2$, Monomial) & 82.00 & 59.00 & \textbf{70.50} & 4.6047 & 4.6045 & 389.31 \\
\midrule
\textbf{ChebyMerge ($d=2$, Ours)} & 81.00 & 67.00 & \textbf{74.00} & 4.6047 & 4.6045 & \textbf{2.75} \\
\textbf{ChebyMerge-CSD ($d=2$, Ours)} & 87.00 & 64.00 & \textbf{75.50} & 4.6047 & 4.6046 & \textbf{2.75} \\
\bottomrule
\end{tabular}
\end{small}
\end{center}
\vskip -0.1in
\end{table*}

Our physical experiments yield several highly significant findings:
\begin{enumerate}
    \item \textbf{The Overfitting-Optimizer Paradox}: Under an unlabeled adaptation stream, minimizing prediction entropy on-the-fly consistently decreases the actual test classification accuracies across all adaptive methods. Because the adaptation stream is small and unlabeled, the optimizer overfits to local transductive noise, driving surrogate entropy down while degrading true generalizability. 
    \item \textbf{Monomial Stiffness and Accuracy Collapse}: PolyMerge ($d=2$) suffers from extreme accuracy degradation, collapsing from $81.50\%$ down to $70.50\%$ (a massive \textbf{-11.00\%} drop). This is due to the extreme ill-conditioning of the monomial basis (condition number $= 389.31$), which creates a stiff, distorted loss landscape where parameter steps are highly uncoordinated and unstable.
    \item \textbf{Chebyshev Landscape Isotropy and Subspace Protection}: In contrast, ChebyMerge ($d=2$) maintains a highly robust, well-conditioned Gram matrix (condition number $= 2.75$, a substantial \textbf{141.8$\times$ improvement}). This well-conditioned, isotropic landscape allows the optimizer to make smooth updates.
    \item \textbf{CSD Controllable Regularization Breakthrough}: By pairing ChebyMerge with Controllable Spectral Decay (CSD), we enforce controllable spectral filtering on top of our well-conditioned workspace. \textbf{ChebyMerge-CSD ($d=2$) achieves the highest accuracy among all continuous subspace adaptation methods, preserving accuracy at $75.50\%$ (outperforming PolyMerge by +5.00\% absolute)}. This provides definitive, physical proof that separating numerical conditioning from regularization is superior: rather than relying on monomial ill-conditioning as an uncontrolled spectral filter, CSD provides explicit, controllable low-pass filtering that successfully shields the network from transductive noise.
\end{enumerate}

\subsubsection{Sensitivity to Learning Rates and Robustness Sweeps}
To thoroughly understand the optimization stability and robustness of each continuous subspace method, we perform a systematic sweep of the base learning rate $\eta \in [10^{-4}, 2 \cdot 10^{-2}]$. Table~\ref{tab:lr_sweep} reports the final test-set classification accuracies for all adaptive merging methods across different learning rates.

\begin{table*}[t]
\caption{Learning Rate Sensitivity and Accuracy Robustness Sweeps}
\label{tab:lr_sweep}
\vskip 0.15in
\begin{center}
\begin{small}
\setlength{\tabcolsep}{3.5pt}
\begin{tabular}{lccccc}
\toprule
Method & $\eta=10^{-4}$ & $\eta=10^{-3}$ & $\eta=5\cdot 10^{-3}$ & $\eta=10^{-2}$ & $\eta=2\cdot 10^{-2}$ \\
\midrule
AdaMerging & \textbf{82.00} & 81.00 & 81.00 & 78.00 & 76.00 \\
PolyMerge ($d=2$) & 81.00 & 80.50 & 75.50 & 70.50 & 66.00 \\
ChebyMerge ($d=2$) & 80.50 & 80.00 & 75.00 & 74.00 & 71.00 \\
ChebyMerge-CSD ($d=2$) & 81.50 & \textbf{80.50} & \textbf{78.00} & \textbf{75.50} & \textbf{70.00} \\
\bottomrule
\end{tabular}
\end{small}
\end{center}
\vskip -0.1in
\end{table*}

This sweep yields critical insights:
\begin{enumerate}
    \item \textbf{High-Rate Catastrophic Collapse of Monomials}: As the learning rate increases, PolyMerge's accuracy collapses catastrophically, dropping from $81.00\%$ at $\eta=10^{-4}$ down to an extremely poor $66.00\%$ at $\eta=2\cdot 10^{-2}$ (a massive \textbf{-15.00\%} collapse). Because the monomial basis is highly ill-conditioned (condition number $= 389.31$), larger gradient steps along distorted dimensions cause the optimization to diverge rapidly.
    \item \textbf{Chebyshev Robustness and Graceful Degradation}: In contrast, ChebyMerge exhibits exceptionally graceful degradation under high learning rates, maintaining $71.00\%$ even at the extreme rate of $\eta=2\cdot 10^{-2}$ (outperforming PolyMerge by \textbf{+5.00\%} absolute). This confirms that ChebyMerge's well-conditioned, isotropic landscape provides a critical safety buffer, preventing catastrophic optimizer divergence.
    \item \textbf{Regularization Leadership of CSD}: Across almost all learning rates, ChebyMerge-CSD consistently achieves the highest accuracy among all continuous subspace methods, illustrating that explicit controllable low-pass filtering is far more robust and stable than relying on monomial ill-conditioning for damping.
\end{enumerate}

\subsection{Methodological Limitations and Architectural Scaling}
\label{sec:limitations}

While ChebyMerge offers a mathematically elegant and numerically superior framework for continuous subspace model merging, intellectual honesty demands a thorough examination of its methodological limits, robustness properties, and topological assumptions.

\subsubsection{Real-World Scope and Open Challenges of the Physical Validation}
We explicitly acknowledge that while our physical evaluation on pre-trained \texttt{CLIP-ViT-B/32} successfully validates the \emph{accuracies, optimization, and conditioning} advantages of ChebyMerge on actual weight tensors and real image datasets, it represents a controlled research setup. Specifically:
\begin{enumerate}
    \item \textbf{Scale of the Adaptation Stream}: Our test-time adaptation stream consists of 100 images, and test evaluation is done on 200 images. While sufficient to cleanly expose the Overfitting-Optimizer Paradox and demonstrate the optimization advantages of Chebyshev subspaces, real-world deployment in large-scale pipelines will involve thousands of samples.
    \item \textbf{Task Dimensionality}: We evaluate merging of $K=2$ tasks (MNIST and SVHN) which represents a standard sequential merging benchmark, but merging a larger number of tasks (e.g., $K \ge 8$ as in the task vectors benchmark) remains an open challenge for continuous parameterizations.
\end{enumerate}
Nevertheless, the successful completion of this experiment on actual pre-trained and fine-tuned foundation models provides definitive, empirical validation of our method on actual deep neural network weights. Scaling this to larger, ultra-deep Language Models (such as LLaMA-3-8B) is a key direction for our future work.

\subsubsection{Robustness of the CSD Decay Factor}
The Controllable Spectral Decay (CSD) framework introduces the hyperparameter $\gamma_{\text{CSD}} \in (0, 1]$ to scale the learning rates of higher-order Chebyshev coefficients. Because test-time adaptation is entirely unsupervised and operates on unlabeled local streams, tuning $\gamma_{\text{CSD}}$ on-the-fly via cross-validation is impossible. However, CSD inherits a strong structural prior from signal processing, where decay rates are chosen based on the expected noise frequency characteristics.

In our empirical robustness sweeps (analyzed in Section~\ref{sec:csd_robustness} of the Appendix), we find that performance is highly robust to the precise selection of $\gamma_{\text{CSD}}$. Any value in the range of $\gamma_{\text{CSD}} \in [0.5, 0.8]$ effectively dampens high-frequency transductive noise without suppressing the lower-degree underlying sensitivity profile. So long as the higher-degree coefficients are not completely frozen ($\gamma_{\text{CSD}} = 0$) or entirely unregularized ($\gamma_{\text{CSD}} = 1.0$), ChebyMerge stably out-generalizes unconstrained methods, making CSD highly robust and practical in blind test-time settings.

\subsubsection{Sequential Depth vs. Branched Topologies}
Mapping the layer index $l$ to a 1D continuous coordinate $x_l \in [-1, 1]$ assumes a strictly sequential, single-path model topology. While sequential structures are standard in many Vision Transformers, modern deep architectures often incorporate parallel residual branches, multi-path layers, or Mixture-of-Experts (MoE) routing. 

As detailed in Section~\ref{sec:complex_topologies} of the Appendix, ChebyMerge can handle multi-path architectures by performing a topological sort to linearize the graph into a sequential depth sequence before mapping to $x_l$. However, we acknowledge that a strictly 1D sequential projection is an approximation for highly branched models. A natural mathematical extension of our work is to formulate multi-dimensional or graph-spectral Chebyshev projections, mapping coefficients directly on the network's topological graph spectrum.

\subsubsection{Asymmetric Sensitivity Profiles}
ChebyMerge's foveated sensitivity matching (implicit boundary concentration) assumes a symmetric profile where the earliest and deepest layers are the most sensitive. While this symmetric profile is widely documented in standard Transformer and CNN architectures \cite{zhang2019all}, certain networks may exhibit highly asymmetric sensitivity (e.g., concentrated exclusively in deep layers) or localized intermediate sensitivity.

In such cases, the symmetric boundary concentration of the standard Chebyshev grid could theoretically be sub-optimal. Crucially, ChebyMerge can easily accommodate asymmetric sensitivity by applying a coordinate-warping diffeomorphism (such as a Beta cumulative distribution function) to the normalized layer index before evaluating the Chebyshev polynomials:
\begin{equation}
\tilde{x}_l = 2 \cdot I_{\bar{l}}(a, b) - 1
\end{equation}
where $I_{\bar{l}}(a, b)$ is the regularized incomplete beta function. This warping allows researchers to custom-shift the representation resolution and concentration of the Chebyshev basis to match any known physical layout, maintaining perfect numerical conditioning while providing complete architectural flexibility.

